\section{Threat Model}
\label{sec:threat-model}
% CLAUDE.md threat model: honeypot-ML integrations retrain the detector on decoy traffic,
% auto-labelled malicious by policy ("everything the honeypot sees is malicious"), as free ground
% truth. The unexamined assumption is that the label is trustworthy because the source is a
% decoy; an attacker who knows or can infer that policy controls exactly what the decoy sees, so
% the honeypot is a write channel into the defender's training set that the defender opens
% deliberately. The poison budget is not a fraction the attacker must fight for -- it is the
% entire honeypot-derived share of the training set that the attacker's traffic is admitted into.
% Narrowed per external review: the labelling policy need not be literally public for the attack
% to work, only known or reliably inferable by the attacker -- see "Attacker knowledge" below.
The systems we attack retrain on two sources: a trusted seed set the defender controls, and
whatever traffic reaches a honeypot, auto-labelled malicious by policy — anything a decoy sees is
unsolicited by definition, so the argument goes, and therefore free ground-truth-malicious data.
That argument holds only if the label is trustworthy because the source is a decoy. It is not,
whenever the attacker knows or can reliably infer that policy: it can then choose exactly what
the decoy sees, with a label it knows the defender will trust. The honeypot becomes a write
channel into the defender's training set, opened by the defender's own design. The poison budget
is not a fraction the attacker has to fight for against a defended channel; it is the entire
honeypot-derived share of the training set that the attacker's admitted traffic reaches, by
construction of the design — what fraction of \emph{all} traffic entering retraining that share
amounts to is a separate, measured quantity (realized poison ratio, below).

% DECISIONS.md s4: attacker capability. Can send arbitrary traffic to the decoy; the pool is
% two-class (benign-shaped and attack-shaped) so A1 (point the benign generator at the decoy)
% and S0/A4 can both be expressed from the same adversary interface; sized (~20k rows) so
% poison ratios up to 50% are reachable with FRESH rows, not resampled duplicates -- a
% resampled duplicate is closer to upweighting an existing sample than to injecting new poison.
\textbf{Attacker capability.} The attacker can send arbitrary traffic to the honeypot, but cannot
touch \code{seed_train}, \code{trusted_eval}, or the model directly — every effect it has
on the deployed detector routes through what the honeypot logs. The pool it draws from is
two-class, sized so poison ratios up to 50\% of the training set are reachable with fresh rows
rather than resampled duplicates; for a tree model, a resampled duplicate is closer to
upweighting a sample already present than to injecting new poison, which would understate what
the attack can do.

% Added per external review: knowledge was an unstated assumption; make each axis explicit and
% say plainly which level A1 as implemented needs, rather than leaving "public" to do that work.
\textbf{Attacker knowledge.} We separate what the attacker needs to know along four axes, rather
than bundling them into one assumption.
\emph{Labelling-policy knowledge} is the one the attack depends on most, and we model it at four
levels: \emph{unknown} (the attacker has no information about how honeypot traffic is used);
\emph{inferred} (the attacker suspects or deduces, from observing honeypot behaviour or general
knowledge of how such systems are built, that anything the decoy logs is treated as malicious
ground truth); \emph{known} (the attacker knows this outright, because the design is published —
an academic description, vendor documentation, an open-source implementation — or because a
honeypot that must be reachable to do its job cannot keep that much secret); and \emph{known +
retraining-aware} (the attacker additionally knows or tracks the defender's retraining schedule
well enough to time its traffic). A1 as implemented assumes the \emph{known} level: it requires
that the attacker knows, not merely suspects, that honeypot traffic is auto-labelled malicious
and fed into training, but it needs nothing beyond that — no retraining-schedule awareness, no
behavioural inference from a black-box honeypot. This is a real assumption, not a vacuous one: an
attacker who only observes the honeypot's external behaviour, without access to a description of
the design, would have to infer the policy rather than know it, and we do not evaluate how much
harder that makes the attack.
\emph{Model knowledge}: none. A1 never queries the detector, computes a gradient against it, or
otherwise interacts with it; everything it needs routes through the honeypot.
\emph{Benign-distribution knowledge}: A1 as implemented draws its poison from benign \emph{feature
rows} already in the attacker's reach in our experiments — the same pool the defender's own
traffic generator uses — which assumes access to the benign class's statistics in the detector's
own feature space. Realizing equivalent mimicry from packet-level traffic alone, without that
feature-space access, is not demonstrated here and is left as future work
(Section~\ref{sec:discussion}).
\emph{Retraining-schedule knowledge}: none needed. The attack is agnostic to when retraining
happens; it only needs its traffic admitted into whatever batch is next retrained on.

% CLAUDE.md scenarios table.
\textbf{Attacker goal.} Raise the detector's false-positive rate on production traffic, or
suppress its detection of a real attack. \textbf{A1 (benign mimicry)}, this paper's subject,
does the first: point the same kind of traffic generator that produces legitimate benign flows
at the honeypot instead of production, and let the auto-labeller stamp the result malicious.
\textbf{S0 (clean)} — the honeypot sees genuine attacks, correctly labelled — is not an attack;
it is the sanity check this design has to pass before any attack against it means anything
(Section~\ref{sec:attack}). \textbf{A2 (boundary drift)} and \textbf{A4 (fingerprint-and-split)}
are scoped out: A2 needs a gradual multi-round adversary this paper does not build, and A4's
controlled measurement arm cannot occur naturally on CICIDS2017, where every family runs on a
single day (Section~\ref{sec:methodology}).

% Added per external review: "realized poison ratio" was used throughout without an operational
% definition.
\textbf{Realized poison ratio.} Where this paper reports a poison ratio, it means the operational
quantity: attacker-admitted honeypot samples divided by all samples entering that round's
retraining (seed plus honeypot, admitted after any defense hook). It is realized, not nominal —
the target ratios named in Section~\ref{sec:methodology} are generation targets that the loop's
batch sizing is tuned to reach by the final round, not a guarantee independent of what else is
happening in the honeypot. In particular this ratio depends on two things we do not model: how
much background (non-adversarial) traffic the honeypot also receives, which would dilute the
attacker's share of what enters training, and any rate limiting the honeypot or its ingestion
pipeline applies, which would cap how much poison the attacker can inject per round regardless of
how much traffic it sends. Every ratio reported in this paper is the realized value measured
directly from the loop's own admitted batches, not back-computed from a generation target, but
it is realized against a honeypot with no competing traffic and no rate limit — a best case for
the attacker that a deployed honeypot need not offer.

% DECISIONS.md s18 "Realized fidelity, not the knob": the capability axis is a measured
% distance, not an assumed one. s26.2/26.3: median cannot resolve the FPR cliff; twin fraction,
% a tail statistic, does.
\textbf{Capability axis.} A1's only free parameter is how closely the attacker's traffic
resembles genuine production benign traffic. We do not assume this from a generation setting;
we measure it directly as the realized nearest-neighbour distance from the injected poison to
trusted benign traffic (Section~\ref{sec:methodology}), and report damage as a function of that
measured distance rather than of the knob that produced it. The median of this distance turns
out to be too coarse a summary: it changes by about 1\% across the exact step where the FPR
channel's damage collapses from catastrophic to noise-floor (Section~\ref{sec:attack}). We
additionally report \emph{twin fraction} — the share of a poison batch whose nearest trusted-
benign neighbour sits within a fixed RMS delta (0.0015, chosen from where the distribution's
lower tail moves between the jitter steps that separate damaging from non-damaging poison,
Section~\ref{sec:attack}) — which does track the collapse and is the primary axis for the FPR
figure; median distance is kept as a secondary statistic since it still separates points once
twin fraction has gone uniformly to zero.
